\documentclass[11pt]{article}
\usepackage[margin=0.9in]{geometry}
\usepackage[T1]{fontenc}
\usepackage{lmodern,amsmath,amssymb,amsthm,booktabs,array}
\usepackage[colorlinks=true,linkcolor=blue,citecolor=blue,urlcolor=blue]{hyperref}
\newtheorem{theorem}{Theorem}
\newtheorem{corollary}[theorem]{Corollary}
\newtheorem{lemma}[theorem]{Lemma}
\newcommand{\N}{\mathbb N}
\newcommand{\gen}[1]{\langle #1\rangle}
\title{V\'elez Caicedo and Cadavid's numerical-semigroup realizability question:\\prescribed maxima and reduced spectral data}
\author{Haobo Ma\\\small ChronoAI Pte Ltd, Singapore
\and Wenlin Zhang\\\small National University of Singapore\\\small ChronoAI Pte Ltd, Singapore}
\date{November 18, 2026}
\begin{document}
\maketitle
\begin{abstract}
We answer the realizability question of V\'elez Caicedo and Cadavid
(\cite[\S5, p.~17, second paragraph]{VC}, arXiv:2611.14316v1) by reflecting
semigroup generators about a prescribed maximum jump: every numerical
semigroup $T$ is realizable, and a positive integer $M$ is an admissible maximum
exactly when $M\in T$ and every minimal generator $h$ satisfies
$h<2M$ and $2M-h\in T$.
The construction gives the unique inclusion-minimal jump set at each
admissible maximum and realizes every sufficiently large maximum.
Using the source's spectral formulas, we also recover the semigroup from
two global dimensions and $M-2$ restricted dimensions when $M\ge2$.
For prime maxima, eventual dimensions at proper-divisor moduli alone need
not distinguish semigroups.
\end{abstract}

\section{The question and its arithmetic form}
We use $\N=\{0,1,2,\ldots\}$. A numerical semigroup is an additive submonoid
$T\subseteq\N$ with finite complement, including $T=\N$.
For $A\subseteq\N$, $\gen A$ is the monoid of finite nonnegative integer
linear combinations of elements of $A$, including the empty sum $0$.
These are the conventions of \cite{RG}; all arithmetic facts needed below
are proved here.

For a finite nonempty positive jump set $S$, put $M=\max S$ and
\[
 \Sigma_S=\gen{\{M\}\cup\{M-s:s\in S\}\cup\{M+s:s\in S\}}.
\]
V\'elez Caicedo and Cadavid study this semigroup through the filtered
spectral dimensions of circulant graphs of orders $n>2M$.
Their final discussion asks which numerical semigroups occur, whether
fewer congruence restrictions suffice to distinguish them, and which
spectral data reconstruction requires \cite[\S5, p.~17]{VC}.
We call $M$ \emph{admissible for $T$} if $\Sigma_S=T$ for some such $S$
with $\gcd S=1$ and maximum $M$.
The source leaves $M$ free. The prescribed-maximum criterion below
strengthens its existential question; the reconstruction consequences use
its already proved dimension formulas.
The quotient-genus identities of Adeniran et al.\ \cite[Theorem 1.3]{AB}
concern quotients of a given numerical semigroup. Here the arithmetic
construction instead starts from an arbitrary target semigroup and
prescribes the maximum of a realizing jump set.

Let $c$ be the conductor of $T$, the least nonnegative integer for which
$c+\N\subseteq T$, and let $m$ be its least positive member.
Thus $c=0$ exactly when $T=\N$.
A positive element of $T$ is \emph{indecomposable} if it is not a sum of
two positive elements of $T$.

\begin{lemma}\label{lem:generators}
The indecomposable elements form the unique minimal positive generating
set $H$ of $T$. This set is finite, nonempty, has $\gcd H=1$, and
\[
 \max H\le\max\{m,c+m-1\}.
\]
\end{lemma}
\begin{proof}
Every $x\in T$ with $x\ge c+m$ and $x>m$ decomposes as
$m+(x-m)$, with $x-m\in T$ positive. Thus an indecomposable is either
$m$ or is smaller than $c+m$, proving the bound and finiteness.
The element $m$ is indecomposable. Strong induction on a positive
$x\in T$ expresses it using indecomposables: either it is one, or it
splits into two smaller positive elements to which induction applies.
Conversely every positive generating set must contain each indecomposable,
since an expression for it can have only one positive summand.
This proves minimality and uniqueness.
Finally, a divisor common to $H$ divides all of $T$, hence divides the
consecutive positive members $c+1$ and $c+2$, and is one.
\end{proof}

\section{Prescribed maxima and the smallest jump set}
\begin{theorem}\label{thm:criterion}
Let $H$ be the minimal positive generating set of a numerical semigroup
$T$. A positive integer $M$ is admissible for $T$ if and only if
\begin{equation}\label{eq:criterion}
 M\in T,\qquad h<2M\ \hbox{and}\ 2M-h\in T\quad(h\in H).
\end{equation}
When these conditions hold, the set
\begin{equation}\label{eq:minimal}
 S_{\min}(T,M)=\{M\}\cup\{|M-h|:h\in H,\ h\ne M\}
\end{equation}
realizes $T$, has gcd one, and is contained in every jump set realizing
$T$ with maximum $M$. It is therefore the unique inclusion-minimal such
set and has the smallest possible number of jumps.
\end{theorem}
\begin{proof}
Suppose $\Sigma_S=T$ with maximum $M$. The defining generators of
$\Sigma_S$ belong to $T$, in particular $M\in T$.
An indecomposable $h\in H$ must itself be one of the nonzero defining
generators: an expression with two positive summands would decompose it.
It cannot be $2M$, which is $M+M$. Thus either $h=M$, or
$h=M-s$ or $h=M+s$ for some $s\in S$ with $0<s<M$.
In the latter cases the paired generator $2M-h=M+s$ or $M-s$
also belongs to $T$, and $h<2M$. For $h=M$ the same conditions hold.
Moreover $|M-h|=s\in S$ whenever $h\ne M$. Since also $M\in S$,
this proves both necessity and the asserted containment.

Conversely assume \eqref{eq:criterion}. Each jump $|M-h|$ with $h\ne M$
is positive and smaller than $M$, so \eqref{eq:minimal} has maximum $M$.
Its defect generators are $M,0,2M$ and the pairs
$\{h,2M-h\}$ for $h\in H\setminus\{M\}$.
They all belong to $T$, and they contain every member of $H$,
with $h=M$ supplied by $M$ itself. Their generated monoid is therefore $T$.
If a positive integer $d$ divides all jumps, it divides $M$ and
each $|M-h|$, hence every $h\in H$; Lemma~\ref{lem:generators} gives $d=1$.
The already proved containment gives minimality and uniqueness.
\end{proof}

\begin{corollary}[All sufficiently large maxima]\label{cor:large}
For any finite nonempty positive generating set $G$ of $T$, every
$M\ge\max\{c,\max G\}$ is admissible, with realization
\[
 S_M=\{M\}\cup\{M-h:h\in G,\ h<M\},\qquad
 |S_M|=1+|G|-\mathbf1_{\{M\in G\}}.
\]
In particular every numerical semigroup is realizable, and every
$M\ge\max\{1,c+m-1\}$ is admissible. The least admissible maximum
can be found by testing \eqref{eq:criterion} for
$1\le M\le\max\{1,c+m-1\}$.
\end{corollary}
\begin{proof}
Here $M\in T$. For $h<M$ the paired generator $2M-h$ is larger than
$M\ge c$, hence in $T$. The defect generators of $S_M$ lie in $T$ and
contain $G$, proving equality; the gcd argument in the theorem applies
because every positive generating set has gcd one by the final argument
of Lemma~\ref{lem:generators}. Reflection is injective and never gives $M$,
which proves the count. For the final bound use $G=H$ and the lemma.
If $c>0$, then $c\ge1$ and $c+m-1\ge m,c$; if $c=0$, then
$T=\N$, $m=1$, $H=\{1\}$ and $S=\{1\}$ works at $M=1$.
\end{proof}

\begin{corollary}[Ordinary semigroups]\label{cor:ordinary}
For $T_m=\{0\}\cup\{m,m+1,\ldots\}$, $m\ge1$, the admissible maxima
are exactly
\[
 M\ge \left\lceil\frac{3m-1}{2}\right\rceil.
\]
\end{corollary}
\begin{proof}
The minimal generators are $m,\ldots,2m-1$. An admissible $M$ belongs
to $T_m$, so $M\ge m$ and every $2M-h$ in the criterion is positive.
It belongs to $T_m$ exactly when it is at least $m$.
The strongest of these inequalities is
$2M-(2m-1)\ge m$, equivalent to the displayed bound.
Conversely this bound implies $M\ge m$, $h<2M$ and $2M-h\ge m$
for every minimal generator, so Theorem~\ref{thm:criterion} applies.
\end{proof}

For example, $T=\gen{3,5}$ has conductor $8$ and least admissible
maximum $8$, where $S_{\min}=\{3,5,8\}$. To check minimality, the only
positive members of $T$ below $8$ are $3,5,6$: the criterion fails on
$h=5,3,5$, respectively, whose partners are $1,7,7\notin T$.
At $M=8$ both partners, $13$ and $11$, lie in $T$.
For $T=\N$, every $M\ge1$ works: the smallest jump set is $\{1\}$
at $M=1$ and $\{M-1,M\}$ otherwise.

\section{Fewer eventual spectral dimensions}
For clarity we retain the source's spectral conventions.
Let $A_{n,S}$ be the adjacency matrix of the undirected circulant graph
with jumps $\{\pm s:s\in S\}$ on $\mathbb Z/n\mathbb Z$, $n>2M$,
and put $r=|S|$ and
$G_h(n)=2^{-h}\operatorname{tr}(A_{n,S}^h)-r^h$ for $h\ge1$.
The algebra $\mathcal A_S=\mathbb Q[G_1,G_2,\ldots]$ is filtered by assigning
weight $h$ to $G_h$. Write $\rho^{\mathrm{glob}}_S(D)$ for the dimension
of its weight-at-most-$D$ part and $\rho_{S;q,a}(D)$ for the dimension
after restricting graph orders to $n\equiv a\pmod q$.
For $T=\Sigma_S$, let $g$ be the number of gaps and
$g_t=|\{v\in\N\setminus T:v\equiv t\pmod M\}|$.
The source proves that, at a common sufficiently large degree $D$,
\begin{align}
 \rho^{\mathrm{glob}}_S(D)&=MD+\lfloor D/2\rfloor+1-M(M+1)-g,
       \label{eq:global}\\
 \rho_{S;M,-t\bmod M}(D)&=D+\lfloor D/2\rfloor-1-g_t
       \quad(0\le t<M).
       \label{eq:restricted}
\end{align}
These are cited formulas, not additional spectral theorems proved here
\cite[Theorems 1.1--1.3, Corollaries 3.9--3.10]{VC}.

\begin{corollary}\label{cor:spectral}
If $M\ge2$, two global dimensions at $D,D+2$ and the $M-2$ dimensions
restricted to $n\equiv-t\pmod M$, $1\le t\le M-2$, at degree $D$
determine $M$ and $\Sigma_S$. All degrees must lie in the source's
common stable range.
If $M=1$, then $\Sigma_S=\N$.
\end{corollary}
\begin{proof}
First recover $M=(\rho^{\mathrm{glob}}_S(D+2)-\rho^{\mathrm{glob}}_S(D)-1)/2$
and then $g$ from \eqref{eq:global}.
Since $M\in T$, its nonnegative multiples belong to $T$, so $g_0=0$.
Equation~\eqref{eq:restricted} recovers $g_1,\ldots,g_{M-2}$ and
$g_{M-1}=g-\sum_{t=1}^{M-2}g_t$ supplies the remaining count.
Let $w_t$ be the least member of $T$ congruent to $t$ modulo $M$.
Such a member exists by cofiniteness. Closure under addition of $M$
shows that the gaps in this class are precisely
$t,t+M,\ldots,w_t-M$, whence $w_t=t+Mg_t$.
Thus $T=\bigcup_{t=0}^{M-1}(w_t+M\N)$ is determined.
For $M=2$ the list of restricted values is empty and the same argument
applies. For $M=1$, $1\in T$ gives $T=\N$.
\end{proof}

\begin{corollary}\label{cor:obstruction}
For every prime maximum $M\ge5$, eventual dimensions at all moduli
properly dividing $M$, even together with the eventual global dimensions,
need not distinguish defect semigroups.
\end{corollary}
\begin{proof}
Take $T_1=\N\setminus\{1,2\}=\gen{3,4,5}$ and
$T_2=\N\setminus\{1,3\}=\gen{2,5}$.
Their conductors are $3,4$, respectively. Both are realizable at every
$M\ge5$ by Corollary~\ref{cor:large}, and both have genus $2$.
The only positive proper divisor of prime $M$ is $1$, for which the
restriction is the unrestricted sequence. Formula~\eqref{eq:global}
therefore makes all these eventual dimension data identical, although
$T_1\ne T_2$.
\end{proof}
This obstruction concerns eventual dimensions at divisor moduli, not
the complete low-degree filtration or arbitrary nondivisor moduli.
Corollary~\ref{cor:spectral} is a sufficient smaller family, not a proof
that $M-2$ restrictions are always necessary.

\section{Coverage and verification}
The source poses questions rather than an assertion of nonrealizability.
The following separates its questions from the additional consequences.
\begin{center}
\small
\begin{tabular}{@{}p{0.43\linewidth}p{0.49\linewidth}@{}}
\toprule
Source item in \cite{VC} & Coverage\\
\midrule
\S5, p.~17, paragraph 2: realizability & Answered: yes, for every numerical
semigroup (Corollary~\ref{cor:large}).\\[3pt]
Same paragraph: smaller sufficient congruence families & Answered: yes
for eventual data, using two global values and $M-2$ restricted values
(Corollary~\ref{cor:spectral}).\\[3pt]
Same paragraph: necessary reconstruction data & Partial: proper-divisor
moduli can fail (Corollary~\ref{cor:obstruction}); a full minimal-data
classification is not covered.\\[3pt]
Prescribed maxima and smallest jump sets & Not posed; proved in
Theorem~\ref{thm:criterion}.\\[3pt]
Ordinary semigroups' least maximum & Not posed; proved in
Corollary~\ref{cor:ordinary}.\\[3pt]
Corollary 3.10: jump-set nonuniqueness & Prior result cited; the unique
minimal set here is for prescribed $T,M$, not uniqueness of every realization.\\
\bottomrule
\end{tabular}
\end{center}

The existential realization statement, including $T=\N$, is the scope of
\texttt{D5.S3.Arith.Additive.}\allowbreak
\texttt{NumericalSemigroupRealizability.result} in the
\href{https://github.com/the-omega-institute/trureturing/blob/8acf7efa41af7fc9ddbbd39571e7ef73462f4509/D5/S3/Arith/Additive/NumericalSemigroupRealizability.lean}{commit-pinned Lean file}.
For a cofinite additive submonoid of $\N$, this declaration gives a
positive $M$ and a finite nonempty positive set $S$ with maximum $M$,
gcd one, and generated monoid equal to $T$.
The prescribed-maximum criterion, minimality, quantitative bounds and
spectral consequences have the written proofs above; the spectral input
is precisely \eqref{eq:global}--\eqref{eq:restricted} from \cite{VC}.

The accompanying \texttt{verify.py} checks exact finite instances of the
construction, fixed-maximum criterion, minimality and residue-count
reconstruction. It enumerates numerical semigroups by their finite gap
sets and also enumerates positive primitive jump sets at bounded maxima.
Equality is compared using complete least-residue representative lists,
not just an initial segment. Exact ranges and outputs are in
\texttt{verification.json}; these finite calculations complement the
general written proofs.

\begin{thebibliography}{9}
\bibitem{AB}
A.~Adeniran, S.~Butler, C.~Defant, Y.~Gao, P.~E.~Harris, C.~Hettle,
Q.~Liang, H.~Nam and A.~Volk,
\emph{On the genus of a quotient of a numerical semigroup},
Semigroup Forum \textbf{98} (2019), 690--700;
arXiv:1809.09360v2, 24 November 2018.
\url{https://doi.org/10.1007/s00233-018-9989-3}.
\bibitem{VC}
J.~D. V\'elez Caicedo and C.~A. Cadavid,
\emph{Reconstructing Numerical Semigroups from Circulant Spectral Dimensions},
arXiv:2611.14316v1, 10 November 2026.
\url{https://arxiv.org/abs/2611.14316v1}.
\bibitem{RG}
J.~C. Rosales and P.~A. Garc\'ia-S\'anchez,
\emph{Numerical Semigroups}, Developments in Mathematics, vol.~20,
Springer, 2009.
\url{https://doi.org/10.1007/978-1-4419-0160-6}.
\end{thebibliography}
\medskip\noindent\small
\copyright\ 2026 Haobo Ma and Wenlin Zhang.
Licensed under Creative Commons Attribution 4.0 International (CC BY 4.0).
\end{document}
